\begin{lemma}\label{lem:decic236-cuspidal-pencil}
Every degree-$10$ morphism
$X_1(2,36)_{\mathbb F_5}\to\mathbb P^1$
has a rational fibre whose support is not entirely cuspidal.
\end{lemma}
\begin{proof}
The source has $24$ rational points, all cuspidal, and its cusp
signature is $1^{24},2^8,6^4$. Suppose that all six rational
fibres are cuspidal. Let $n_1\geq\cdots\geq n_6$ count the
distinct rational cusps in these fibres. Then
\[
 \sum_i n_i=24,\qquad n_1\geq4,\qquad
 n_2\geq\left\lceil\frac{24-n_1}{5}\right\rceil.
\]
Let $A_1,A_2$ be the two fibres with largest rational support.
For a cuspidal divisor $A=\sum_c m_c c$, use the geometric
coefficient norm
\[
 \|A\|^2=\sum_c \deg(c)m_c^2.
\]
If $\deg A=10$ and it contains $n$ distinct rational cusps,
then $\|A\|^2\leq(11-n)^2+n-1$.
The fibres have disjoint support. Checking $4\leq n_1\leq10$
and the displayed lower bound on $n_2$ gives
\[
 \|A_1-A_2\|^2\leq104.
\]
The positive and negative parts both have degree $10$.

Lift the cusps over a cyclotomic extension unramified at $5$.
Their difference is torsion by Manin--Drinfeld and reduces to a
principal class. Torsion injectivity over this unramified odd local
field makes the lifted difference principal. Its coefficient vector
therefore lies in the full geometric cuspidal modular-unit lattice
and is fixed by Frobenius $5$.

We compute this lattice from the integral period map for
$\Gamma_\Delta(72)$, where $\Delta=\{\pm1,\pm35\}$.
There are $64$ geometric cusps and $36$ Frobenius orbits.
Intersecting the complete principal-divisor lattice with the
Frobenius-fixed cusp-divisor lattice gives a rank-$35$ lattice
in $\mathbb Z^{64}$. Ordinary Euclidean norm in these coordinates
is exactly the weighted norm above.

Complete short-vector enumeration gives $24552$ pairs of nonzero
vectors of squared norm at most $104$. None has positive part of
degree $10$. An independent calculation reconstructs the full
period kernel and obtains the same enumeration using PARI.
This contradicts the vector $A_1-A_2$.
\end{proof}
